\newcommand{\tipoRequerimento}{Solicitação de pagamento de ajuste}
\newcommand{\tese}{CORREÇÃO MONETÁRIA E JUROS DE PAGAMENTO}

\section{DOS FATOS E DOS DIREITOS}

\begin{enumerate}[resume]
    \item O Município, na condição de consumidor dos serviços de distribuição de energia elétrica, identificou a ocorrência de cobrança indevida referente à Unidade Consumidora \unidadeConsumidora{}.   

    \item Diante da constatação da irregularidade, o Município apresentou reclamação administrativa \textbf{REC 004/2025}, em \textbf{21/08/2025}, perante a concessionária requerendo as providências cabíveis para a restituição dos valores cobrados em desacordo com as normas regulamentares.

    \item Após a análise da reclamação, a concessionária reconheceu a cobrança indevida e o direito à restituição. Em decorrência desse reconhecimento, a concessionária realizou o pagamento em favor do Município em \textbf{25/02/2026}, no valor de \textbf{R\$ 17.868,06}.

    \item Ocorre que, após a análise do referido pagamento, o Município constatou que o valor efetivamente restituído não correspondia à integralidade do montante devido, em razão de interpretação equivocada ou erro no cálculo efetuado pela distribuidora.

    \item Diante da insuficiência do ressarcimento, o Município apresentou contestação demonstrando fundamentadamente o equívoco no montante devolvido.

    \item Após a análise do caso, Foi reconhecido o direito do Município à complementação dos valores indevidamente retidos.

    \item Em cumprimento à decisão, a concessionária realizou novo pagamento em favor do Município, de natureza complementar, no valor de \textbf{R\$ 12.799,41}, efetuado em \textbf{11/05/2026}.

    \item Contudo, o pagamento complementar solucionou apenas a parcela principal, não contemplando a devida atualização monetária e os juros de mora incidentes sobre o período em que o valor permaneceu indisponível para o ente público municipal.

    \item Verifica-se, portanto, que entre \textbf{25/02/2026} e \textbf{11/05/2026} transcorreu período de aproximadamente \textbf{3 meses}, durante o qual o Município permaneceu privado da disponibilidade do montante reconhecidamente devido.

    \item Cumpre então proceder com a apuração dos efeitos financeiros decorrentes do período transcorrido, especialmente no tocante à atualização monetária e aos juros de mora aplicáveis.

    \item Dessa forma, permanece controvérsia quanto à recomposição integral dos efeitos financeiros decorrentes do período em que o valor reconhecidamente devido permaneceu sem a devida disponibilização ao Município.

    \item A presente manifestação tem por finalidade submeter novamente a questão à apreciação da ANEEL, a fim de que seja analisada a existência do saldo financeiro decorrente do período de atraso e, sendo confirmado o direito, seja determinada à concessionária a respectiva complementação.
\end{enumerate}

\section{DOS PEDIDOS}

\begin{enumerate}[resume]
\item Diante do exposto, requer-se:
    \begin{enumerate}
    \item O recebimento e a análise da presente manifestação, vinculando-a, se cabível, ao Processo ANEEL nº \textbf{010.619.56026-39} e aos demais processos/protocolos relacionados ao caso;

    \item O reconhecimento de que os pagamentos realizados pela concessionária, no valor de \textbf{R\$ 17.868,06} e \textbf{R\$ 12.799,41}, não contemplaram integralmente a atualização monetária e os juros de mora incidentes sobre o período de indisponibilidade dos valores;

    \item A apuração da atualização monetária e dos juros de mora incidentes sobre o valor devido, considerando o período de \textbf{25/02/2026} e \textbf{11/05/2026} e os critérios regulamentares aplicáveis;

    \item A determinação à concessionária para que recalcule o valor devido, observando os critérios definidos pela ANEEL, e efetue o pagamento do eventual saldo remanescente ao Município;

    \item A apresentação de memória de cálculo detalhada, com a indicação do valor principal, período de incidência, índice de correção, juros aplicados, quantidade de dias e valor final apurado;

    \item A apresentação dos comprovantes de pagamento e dos respectivos cálculos, para conferência do integral cumprimento das obrigações decorrentes da cobrança indevida;

    \item Caso o pedido seja indeferido, que a ANEEL apresente decisão fundamentada, indicando as razões jurídicas para afastar a incidência da atualização monetária e/ou dos juros no período questionado.
\end{enumerate}

\end{enumerate}